% ============================================================================
% Section 4: Agent Coverage Metrics
% ============================================================================

\section{Agent Coverage Metrics}
\label{sec:coverage}

Traditional code coverage metrics---statement, branch, path,
condition---are defined over deterministic control-flow graphs extracted
from source code. AI agents do not have source code in the traditional
sense; their ``code'' is a composition of prompts, tool inventories,
model weights, and orchestration logic. This section defines
agent-specific coverage metrics that measure the thoroughness of a
stochastic test suite.

\subsection{Coverage Dimensions}
\label{sec:coverage:dimensions}

\begin{definition}[Agent Coverage Tuple]
\label{def:coverage-tuple}
The \emph{agent coverage} of a test suite $\mathcal{S}$ executed over
$N$ total trials is a five-dimensional tuple:
\begin{equation}
\label{eq:coverage-tuple}
  \mathcal{C}(\mathcal{S}) = (C_{\text{tool}}, C_{\text{path}},
    C_{\text{state}}, C_{\text{boundary}}, C_{\text{model}})
    \in [0,1]^5
\end{equation}
We define each dimension below.
\end{definition}

\subsubsection{Tool Coverage}
\label{sec:coverage:tool}

\begin{definition}[Tool Coverage]
\label{def:tool-coverage}
Let $\mathcal{T} = \{t_1, \ldots, t_k\}$ be the agent's available tool
set and $\mathcal{T}_{\text{used}} \subseteq \mathcal{T}$ the subset of
tools invoked across all traces in the test suite. \emph{Tool coverage}
is:
\begin{equation}
\label{eq:tool-coverage}
  C_{\text{tool}} = \frac{|\mathcal{T}_{\text{used}}|}{|\mathcal{T}|}
\end{equation}
\end{definition}

Tool coverage is the simplest dimension: it measures whether the test
suite exercises all available tools. An agent with 10 tools but tests
that only trigger 3 has $C_{\text{tool}} = 0.30$, indicating that 7
tools are untested and may harbor latent regressions.

\begin{remark}
Tool coverage does not account for \emph{how} a tool is used---only
whether it is invoked. A tool called with a single trivial input achieves
the same $C_{\text{tool}}$ contribution as one exercised across diverse
inputs. We address input diversity through boundary coverage
(\cref{def:boundary-coverage}).
\end{remark}

\subsubsection{Decision-Path Coverage}
\label{sec:coverage:path}

An agent's execution trace follows a path through a decision space
defined by tool selections, branching conditions in the orchestration
logic, and model outputs at each step.

\begin{definition}[Decision Path]
\label{def:decision-path}
A \emph{decision path} is the sequence of action-tool pairs in a trace:
$\rho(\tau) = ((a_1, t_1), (a_2, t_2), \ldots, (a_m, t_m))$. Two
traces $\tau, \tau'$ have the same decision path if
$\rho(\tau) = \rho(\tau')$.
\end{definition}

\begin{definition}[Decision-Path Coverage]
\label{def:path-coverage}
Let $\mathcal{P}_{\text{obs}}$ be the set of distinct decision paths
observed across all traces and $\mathcal{P}_{\text{est}}$ be an estimate
of the total number of feasible decision paths. \emph{Decision-path
coverage} is:
\begin{equation}
\label{eq:path-coverage}
  C_{\text{path}} = \frac{|\mathcal{P}_{\text{obs}}|}{|\mathcal{P}_{\text{est}}|}
\end{equation}
\end{definition}

The denominator $|\mathcal{P}_{\text{est}}|$ is, in general,
intractable to compute exactly (the path space may be infinite for
open-ended agents). We employ two estimation strategies:

\begin{enumerate}[leftmargin=*]
  \item \textbf{Capture-recapture estimation.} Using the Chao1 estimator
    from ecology: if we observe $d$ distinct paths, of which $f_1$ appear
    exactly once and $f_2$ appear exactly twice, then:
    \begin{equation}
    \label{eq:chao1}
      |\mathcal{P}_{\text{est}}| \approx d + \frac{f_1^2}{2 f_2}
    \end{equation}

  \item \textbf{Orchestration-graph analysis.} For agents with explicit
    orchestration graphs (e.g., LangGraph state machines), enumerate paths
    up to a maximum depth $d_{\max}$, yielding a finite upper bound.
\end{enumerate}

\subsubsection{State-Space Coverage}
\label{sec:coverage:state}

\begin{definition}[Agent State]
\label{def:agent-state}
An \emph{agent state} is a tuple $\sigma = (\ell, \mathbf{v}, h)$ where
$\ell$ is the current \emph{location} in the orchestration graph,
$\mathbf{v}$ is the vector of \emph{state variables} (e.g., accumulated
context, tool outputs, counters), and $h$ is a hash of the
\emph{conversation history} up to a fixed window.
\end{definition}

\begin{definition}[State-Space Coverage]
\label{def:state-coverage}
Let $\mathcal{S}_{\text{obs}}$ be the set of distinct states observed
across all traces. Since the full state space is infinite (due to
natural language content), we project states onto a finite abstraction
$\pi: \Sigma \to \hat{\Sigma}$ and define:
\begin{equation}
\label{eq:state-coverage}
  C_{\text{state}} = 1 - e^{-|\pi(\mathcal{S}_{\text{obs}})| / \lambda}
\end{equation}
where $\lambda > 0$ is a normalization constant set to the expected
number of distinct abstract states for the agent class (calibrated
empirically). This formulation maps to $[0, 1)$ and asymptotically
approaches 1 as more states are discovered.
\end{definition}

The exponential normalization (\cref{eq:state-coverage}) addresses the
open-ended nature of agent state spaces. Unlike tool or path coverage,
which have natural denominators, state-space exploration has no fixed
ceiling. The exponential form ensures diminishing returns: the first 50
distinct states contribute more than the next 50.

\subsubsection{Boundary Coverage}
\label{sec:coverage:boundary}

\begin{definition}[Tool Boundary]
\label{def:boundary}
A \emph{tool boundary} for tool $t_j$ is a pair $(p_j^{\text{name}},
\{v^{\min}, v^{\max}\})$ where $p_j^{\text{name}}$ is a parameter name
and $v^{\min}, v^{\max}$ are the minimum and maximum values of the
parameter's domain (or representative extremes for string/object types).
Let $\mathcal{B}$ denote the set of all tool boundaries.
\end{definition}

\begin{definition}[Boundary Coverage]
\label{def:boundary-coverage}
Let $\mathcal{B}_{\text{tested}} \subseteq \mathcal{B}$ be the set of
boundaries for which at least one trace invoked the tool with a parameter
value at or near the boundary (within tolerance $\epsilon$).
\emph{Boundary coverage} is:
\begin{equation}
\label{eq:boundary-coverage}
  C_{\text{boundary}} = \frac{|\mathcal{B}_{\text{tested}}|}{|\mathcal{B}|}
\end{equation}
\end{definition}

Boundary coverage extends the classical boundary-value analysis to agent
tool invocations. It tests whether the agent exercises edge cases of
its tools---empty lists, maximum-length strings, zero values, negative
numbers---which are common sources of failures.

\subsubsection{Model Coverage}
\label{sec:coverage:model}

\begin{definition}[Model Coverage]
\label{def:model-coverage}
Let $\mathcal{M}_{\text{target}}$ be the set of language models the
agent is intended to support (e.g., GPT-4o, Claude 3.5 Sonnet,
Gemini 2.0 Flash, Llama 3.3 70B) and
$\mathcal{M}_{\text{tested}} \subseteq \mathcal{M}_{\text{target}}$ the
subset tested. \emph{Model coverage} is:
\begin{equation}
\label{eq:model-coverage}
  C_{\text{model}} = \frac{|\mathcal{M}_{\text{tested}}|}{|\mathcal{M}_{\text{target}}|}
\end{equation}
\end{definition}

Model coverage captures a dimension unique to LLM-based systems: the
same agent deployed with different models can exhibit dramatically
different behaviors. An agent tested only with GPT-4o may fail
catastrophically on Claude 3.5 Sonnet due to differences in tool-calling
conventions, reasoning patterns, or output formatting.

\subsection{Aggregate Coverage}
\label{sec:coverage:aggregate}

\begin{definition}[Overall Coverage]
\label{def:overall-coverage}
The \emph{overall coverage} is the geometric mean of the five dimensions:
\begin{equation}
\label{eq:overall-coverage}
  C_{\text{overall}} = \left( \prod_{i=1}^{5} C_i \right)^{1/5}
\end{equation}
\end{definition}

We use the geometric mean rather than the arithmetic mean because it
penalizes imbalance: a test suite with $C_{\text{tool}} = 1.0$ but
$C_{\text{model}} = 0.0$ should receive low overall coverage, as the
zero model coverage represents a critical gap.

\begin{remark}[Weighted Coverage]
In practice, different coverage dimensions may have different importance.
We support weighted geometric means:
$C_{\text{overall}}^{(w)} = \prod C_i^{w_i}$ with
$\sum w_i = 1$, allowing teams to prioritize dimensions relevant to
their deployment context.
\end{remark}

\subsection{Coverage Monotonicity}
\label{sec:coverage:monotonicity}

\begin{theorem}[Coverage Monotonicity]
\label{thm:coverage-mono}
For any test suite $\mathcal{S}$ and any additional test execution:
\begin{enumerate}[leftmargin=*]
  \item $C_{\text{tool}}, C_{\text{boundary}}, C_{\text{model}}$ are
    non-decreasing, and strictly increase when a new tool, boundary,
    or model is observed.
  \item $C_{\text{state}}$ is non-decreasing in the number of distinct
    states observed.
  \item $C_{\text{path}}$ is \emph{asymptotically monotone}: as
    $|\mathcal{P}_{\text{obs}}| \to \infty$, the Chao1 correction
    vanishes and $C_{\text{path}} \to 1$. For finite test suites,
    $C_{\text{path}}$ may temporarily decrease when a new singleton
    path increases the Chao1 denominator faster than the numerator;
    this non-monotonicity is bounded by
    $O(f_1^2 / (f_2 \cdot |\mathcal{P}_{\text{obs}}|^2))$.
\end{enumerate}
\end{theorem}

\begin{proof}[Proof sketch]
Parts (1) and (2) follow directly: numerators are non-decreasing while
denominators are fixed (part~1), or the coverage function is monotone
in its argument (part~2). Part (3) follows from the asymptotic
consistency of the Chao1 estimator~\citep{chao1984nonparametric}: as
$n \to \infty$, $\hat{S}_{\text{Chao1}} \to S_{\text{total}}$ and
$C_{\text{path}} \to 1$.
See \cref{app:proof:coverage} for the formal argument.
\end{proof}

\subsection{Coverage-Guided Test Generation}
\label{sec:coverage:generation}

Coverage metrics naturally drive test generation: the framework
identifies the lowest-coverage dimension and generates inputs designed to
increase it.

\begin{algorithm}[t]
\caption{Coverage-Guided Test Input Generation}
\label{alg:coverage-gen}
\begin{algorithmic}[1]
\REQUIRE Agent $A$, current coverage $\mathcal{C}$, target coverage $\mathcal{C}^*$
\ENSURE New test scenario $S_{\text{new}}$
\STATE $d^* \gets \argmin_{d \in \{1,\ldots,5\}} C_d / C_d^*$
  \COMMENT{Identify weakest dimension}
\IF{$d^* = \text{tool}$}
  \STATE Generate input $x$ designed to trigger unused tools
    $\mathcal{T} \setminus \mathcal{T}_{\text{used}}$
\ELSIF{$d^* = \text{path}$}
  \STATE Generate input $x$ to explore under-represented orchestration
    branches
\ELSIF{$d^* = \text{state}$}
  \STATE Generate input $x$ to reach novel state regions
\ELSIF{$d^* = \text{boundary}$}
  \STATE Generate input $x$ with extreme parameter values for untested
    boundaries
\ELSIF{$d^* = \text{model}$}
  \STATE Select untested model $\mu \in \mathcal{M}_{\text{target}}
    \setminus \mathcal{M}_{\text{tested}}$
\ENDIF
\RETURN $S_{\text{new}} = (x, \mathcal{P}, E)$
\end{algorithmic}
\end{algorithm}

\begin{proposition}[Coverage-Fault Correlation]
\label{prop:coverage-fault}
Higher agent coverage correlates with higher fault detection probability.
Formally, for two test suites $\mathcal{S}_1, \mathcal{S}_2$ with
$\mathcal{C}(\mathcal{S}_1) \geq \mathcal{C}(\mathcal{S}_2)$
(component-wise), the expected number of distinct faults detected
satisfies $\E[F(\mathcal{S}_1)] \geq \E[F(\mathcal{S}_2)]$ under mild
regularity conditions on the fault distribution.
\end{proposition}

This proposition mirrors the well-established correlation between code
coverage and fault detection in traditional software
testing~\citep{zhu1997software}. We validate it empirically in
\cref{sec:experiments:characterization}.
